\chapter{Introduction}
\label{chap:intro}

Firm growth has a natural baseline, and the facts we are after are best read as departures from
it. Gibrat's law of proportionate effect takes the growth of a firm over a period to be a random
shock independent of its current size \citep{gibrat1931}, so that log-size performs a random walk
and firms of every size grow, on average and in variability, alike. Much of the empirical
firm-growth literature is the study of where this baseline breaks \citep{sutton1997}, and it
breaks in two definite ways that we take as our starting point.

The first is in the volatility of the growth rate. Measured year over year,
a firm's growth volatility falls with its size as a power law,
\begin{equation}
  \sigma(S) \propto S^{-\beta},
  \label{eq:intro-sizevar}
\end{equation}
with a small exponent $\beta \approx 0.15$--$0.20$ that is stable across datasets
\citep{stanley1996, amaral1997}. Gibrat's law makes the volatility independent of size,
$\beta = 0$; a firm assembled from many independent parts would instead average them down to
$\beta = 1/2$; the measured exponent lies between the two and matches neither.

The second is in the shape of the growth distribution. Independent multiplicative shocks would
make it Gaussian, but the demeaned growth rate is instead sharply peaked and fat-tailed, close to
the Laplace form \citep{stanley1996, bottazzi2006}, and it keeps that same shape inside size-based
subsets of firms, large firms and small ones alike. We will refer to these two throughout as the
size--variance law and the growth tent. A model of firm growth has to produce both, and the
difficulty is to produce them together.

\section{The granular explanation and where it fails}
\label{sec:intro-granular}

The standard explanation treats a firm as an aggregate of many smaller units, each carrying
its own idiosyncratic shock. Write the firm's size as the
total over its $K$ sub-units, $S = \sum_{i=1}^{K} s_i$. Over a year each sub-unit grows by an
independent random shock, $s_i \to s_i(1 + \eta_i)$, with the $\eta_i$ of mean zero and common
variance $\sigma_0^2$. The firm's growth rate is then the share-weighted average of these
shocks,
\begin{equation}
  g = \frac{\Delta S}{S} = \sum_{i=1}^{K} w_i\, \eta_i,
  \qquad w_i = \frac{s_i}{S},
\end{equation}
and, the shocks being independent, its variance is
\begin{equation}
  \sigma^2(S) = \sigma_0^2 \sum_{i=1}^{K} w_i^2 = \sigma_0^2\, H,
\end{equation}
set entirely by the Herfindahl index $H = \sum_i w_i^2$ of the sub-unit shares.

This already tells us what the size--variance law requires. Suppose first that the units are
of comparable size. Then $w_i \approx 1/K$ and $H \approx 1/K$, and with $K$ proportional to
$S$ we obtain $\sigma(S) \propto S^{-1/2}$, that is $\beta = 1/2$, well above the measured
exponent, so the shares cannot be even. $H$ has to decay more slowly than $1/S$,
and this happens only when a few large units carry a finite share of the firm and their shocks
dominate the sum.

The same mechanism was established first at the level of the whole economy by
\citet{gabaix2011}, with the units being firms and the aggregate GDP. Firm sizes follow a
Zipf-like law \citep{axtell2001}, a heavy tail $P(s > x) \sim x^{-\nu}$ with $\nu$ close to one, and for such a
tail the Herfindahl index does not vanish as $1/N$ but much more slowly,
\begin{equation}
  H \sim N^{-2\left(1 - 1/\nu\right)}, \qquad 1 < \nu < 2,
\end{equation}
because the sum of squared sizes is carried by the largest few. Their shocks do not cancel,
and aggregate fluctuations decay as $N^{-(1 - 1/\nu)}$ rather than $N^{-1/2}$.
\citet{sutton2002} and \citet{wyartbouchaud2003} carried the argument back inside a single
firm. With $K \propto S$ sub-units whose sizes carry the same heavy tail,
$\sigma(S) \propto S^{-(1 - 1/\nu)}$, so
\begin{equation}
  \beta = 1 - \frac{1}{\nu},
  \label{eq:intro-beta-nu}
\end{equation}
which lands anywhere between $0$ and $1/2$ as the tail runs from Zipf ($\nu \to 1$) to thin
($\nu \to 2$). A heavy-tailed distribution of sub-unit sizes thus delivers both facts at once. It
gives a variance that falls too slowly with size, and a growth distribution that stays
fat-tailed because the sum in $g$ is carried by a handful of dominant units. This is the granular
model, and it reproduces both facts.

\citet{moran2024} extract further predictions from the model, and when these predictions are
tested on the data the model fails. The first new prediction comes from dividing each firm's
growth rate by its own volatility. This cancels the firm-to-firm difference, and what remains
is an average of the firm's own independent shocks. By the central limit theorem this average
is closer to Gaussian the more sub-units the firm has, so the largest firms should be the
closest to a Gaussian. The data show the opposite, the rescaled growth staying fat-tailed at
every size, with the largest firms no nearer a Gaussian than the smallest. The second new
prediction comes from dividing each firm's volatility by the typical volatility of firms its
size. This strips out the size--variance trend and leaves only the spread from one firm to the
next, which in the model carries no size dependence of its own, so the distribution of these
ratios should look the same at every size. This part is borne out in the data. What fails is
the tail of that distribution, which in the data is far thinner than the granular model
predicts.

The last new prediction is on the higher moments of the size--variance distribution. Each moment
of that spread falls with size as its own power law, $E[\sigma^q \mid S] \sim S^{-\zeta_q}$, and
the first exponent is the size--variance law again, $\zeta_1 = \beta$. What the test turns on is
how the higher exponents depend on the order $q$. The granular models make a sharp prediction
here. Because a firm assembled from few sub-units keeps a volatility that does not shrink with
its size, the conditional distribution of volatility at a given size carries a heavy tail whose
weight, not whose position, is what varies with size. Every moment beyond the first is then
dominated by that same tail, and the exponents collapse onto a single value independent of the
order, $\zeta_q = \zeta_{q'}$ for all $q, q' \ge 2$. The data refuse this. Their exponents keep
rising with the order, $\zeta_1,\dots,\zeta_4 \approx 0.20,\,0.39,\,0.51,\,0.58$
\citep{moran2024}, a clear dependence on $q$ that the granular composition cannot produce. We
will refer to this $q$-dependence throughout as the multiscaling of the size-conditioned
volatility: the volatility at a given size is governed not by one scale but by a distribution of
them, whose shape changes as the firm grows.

Two reference behaviours bound the effect and will recur. A conditional volatility with a single
scale, a fixed distributional shape whose level alone slides with size, gives exponents locked
to the first, $\zeta_q = q\,\zeta_1$, a straight line in $q$. The granular tail, at the other
extreme, flattens them to a constant beyond $q = 1$. The data lie between the two, rising but
sublinearly; \citet{moran2024} attribute the shortfall below the straight line to subleading
corrections from the incipient tail of their volatility distribution. Whatever its asymptotic
fate, the measured $q$-dependence is the observable that separates the data from the granular
composition, and it is the one this thesis sets out to generate. As to what the data do require,
\citet{moran2024} conclude that the missing ingredient is correlation: that the shocks hitting a
firm become ``more and more correlated as its size increases, due either to supply chain effects
affecting inputs or customers common to all sub-units'', pointing at relations between firms
rather than inside them.

\section{A model of interacting firms}
\label{sec:intro-move}

The sources usually called on for these correlations, shared customers and suppliers, are in
truth links between firms rather than inside one --- the premise of the production-network
account of the macroeconomy, and of the one existing network model of firm volatility
(Chapter~\ref{chap:background}). The granular models, by contrast, place no relation between
firms at all: they are, in the phrase of \citet{moran2024}, ``island models''. We take that
premise one level down, to the firm-growth statistics themselves. We treat the
firms as the interacting parts and the economy as the aggregate, and look for the
size-dependent correlations between companies rather than within a single one.

Concretely, we model the economy as a network of interacting firms, each node a company and
each edge a signed competitive or cooperative coupling, drawn at random into a disordered
interaction matrix. The dynamics are those of the Generalized Lotka--Volterra system from
theoretical ecology, each firm growing logistically and coupled to its neighbours through that
matrix. On its own this system is bounded. Even in the disordered regime where its dynamics stay
chaotic, firm sizes remain bounded rather than growing, so it cannot represent a steadily
expanding economy. The model of this thesis provides one. Its defining change is that the
self-limitation and the
interactions act on each firm's size relative to the average firm rather than on its absolute
size, which makes the dynamics scale-invariant. Scale invariance under a rescaling of all sizes,
and the mean-coupled form it takes, are inherited from the econophysics models of wealth and
firm sizes (Chapter~\ref{chap:background}); what those models leave out, and what we add, is that
the fluctuations are not an injected external noise but are generated by the disordered
interactions themselves, in a persistently fluctuating state from which we read off the growth
statistics of the firms. The model also adds a small immigration term, a
firm-entry rate that keeps firms from reaching exact extinction.
Total output then grows exponentially without diverging, and in the strongly disordered regime
the composition of the economy never settles but keeps changing chaotically. It is in this
state that the two facts appear together, and switching off the disorder removes both, so their
origin lies in the interactions and not in any single firm. The model recovers the
anomalous multiscaling of the higher moments, the test the granular picture fails, and passes
the other two size-conditioned tests of \citet{moran2024} as well; the size--variance exponent
comes out with the right sign as a clean power law, $\beta \approx 0.45$ at our largest sizes
against the empirical $0.15$--$0.20$, closer than before but still not matched in magnitude.
The thesis's sharpest finding is about the mechanism: the size-conditioned statistics are
generated by the degree heterogeneity of the network. Under the standard dilute normalization a
firm's noise variance grows with its connectivity, so the graph's degree distribution is printed
onto the firm-level fluctuations; on homogeneous graphs the multiscaling vanishes exactly, and
only a scale-free network reproduces the empirical curve. The network is not scenery here but
the explanatory ingredient, which is the production-network premise carried down to the
firm-growth statistics.

The rest of the thesis is organized as follows. Chapter~\ref{chap:growing} sets out the model.
It introduces the relative GLV and its scale-invariant growth, says why the obvious GLV will
not do and what the relative interaction fixes, and locates the chaotic, growing regime in
which the firm-growth statistics are measured. Chapter~\ref{chap:firmgrowth} measures the two
facts in the fluctuating state, sets out the estimator that measurement across many
realizations requires, recovers the multiscaling and the other two conditional tests that
separate the data from the granular models, and then dissects the mechanism: swapping the
degree distribution of the graph while holding everything else fixed shows the heterogeneity
to be what generates the size-conditioned statistics. Chapter~\ref{chap:conclusion} draws
the results together and points to what would tie the model more tightly to the data, in
particular a direct measurement of the size-dependent correlations it generates. Two
appendices hold the supporting material: a dynamical mean-field theory of the relative model,
derived in Appendix~\ref{app:dmft}, and the numerical methods
and measurement protocol (Appendix~\ref{app:numerics}).
